\section{Elementary enclosures and alternative gamma construction}
\label{app:elementary}
This appendix records the finite tails behind \cref{lem:elementary}.
For positive rational \(x=2^ky\), \(1\le y<2\), put
\(t=(y-1)/(y+1)\), so \(\abs t\le1/3\). Then
\begin{equation}
 \log y=2\sum_{j=0}^{M-1}\frac{t^{2j+1}}{2j+1}+R_M,
 \qquad\abs{R_M}\le
 \frac{2\abs t^{2M+1}}{(2M+1)(1-t^2)}.
 \label{eq:log-series-provider}
\end{equation}
Use \(t=1/3\) for \(\log2\), with its error multiplied by \(\abs k\).
Positive interval logarithms use monotonicity and separated endpoints.
Square roots use exact square comparisons and bisection after scaling
to \([1,4]\).

For \(\abs t\le1/2\), the first \(M\) arctangent terms have
remainder at most \(\abs t^{2M+1}/(2M+1)\). Reduce
\(1/2\le t\le1\) by
\(\arctan t=\arctan(1/2)+\arctan((t-1/2)/(1+t/2))\),
whose new argument is at most \(1/3\); use oddness and
\(\arctan t=\pi/2-\arctan(1/t)\) for the other ranges.
Tangent addition and the angle range give
\(\pi=4[\arctan(1/2)+\arctan(1/3)]\).
For \(x>0\),
\(\Log(x+\ii y)=\log(x^2+y^2)/2+\ii\arctan(y/x)\).
Rational reduction boundaries can be subdivided; either valid neighboring
formula agrees at the boundary.

Enclose \(a/\log2\) to width at most \(1/8\) and round its rational
midpoint to an integer \(k\). This is an encoding choice. Refinement
with that \(k\) fixed certifies \(\abs{a-k\log2}<2\).
For a reduced real argument of modulus at most two, the degree-\(M\)
Taylor error is at most \(9\,2^{M+1}/(M+1)!\).
For phase reduction choose a nearby multiple of \(2\pi\) from an
enclosure of width \(1/8\); its residual has modulus at most
\(9\pi/8<4\). The corresponding complex Taylor error is at most
\(81\,4^{M+1}/(M+1)!\).
Refine \(\pi\) and \(\log2\) with the actual chosen multipliers.
Relative real-exponential error \(\tau\) is obtained by absolute
Taylor error at most \(\tau/9\), since the reduced exponential
is at least \(1/9\), together with a separate input-error allowance.
For complex input perturbation use
\[
 \abs{\e^{z+q}-\e^z}\le\e^{\Rea z}(\e^{\abs q}-1).
\]
If \(\abs a\ge d_0>0\) and \(\abs{\widehat a-a}\le\eta<d_0\),
then the reciprocal error is at most \(\eta/[d_0(d_0-\eta)]\).

Every truncation is chosen by a strict rational tail test. For a reduced
series with \(b\)-bit rational input and requested \(P\)-bit accuracy,
\(O(P+b)\) terms suffice. Explicit common denominators such as
\(q^{2M+1}(2M+1)!\) have \(O((P+b+10)^2)\) bits.
Elementary integer arithmetic gives the conservative
\(O((P+b+10)^6)\) work bound used in \cref{thm:bit-cost}.
This refers to reduced arguments and mantissa accuracy. Absolute
accuracy after rescaling an exponential additionally counts its output
magnitude; there is no polynomial absolute-error assertion for an
unrestricted enormous exponential.

\subsection{Euler-product fallback and its internal error}
\label{app:euler-provider}
The elementary construction of \cite{rec09} remains useful as a second
finite formula. Let \(\zeta=(1+\ii u)/4\), \(v=\zeta-1\),
choose rational \(B_\Gamma\ge\abs v\), and take integer
\(s\ge\max(2B_\Gamma,2)\). Set \(q=B_\Gamma/(s+1)\le1/2\) and
\[
 F_s(v)=v\log s-\sum_{j=1}^s\Log(1+v/j).
\]
Euler's product \cite{dlmfgammaproduct} and the harmonic-sum remainder
give
\begin{align}
 \abs{F_s-\log\Gamma(1+v)}
   &\le B_\Gamma/s+B_\Gamma^2/[2s(1-q)],\nonumber\\
 \abs{F_s'-(\log\Gamma)'(1+v)}
   &\le1/s+B_\Gamma/[s(1-q)],\nonumber\\
 \left|\frac{F_s^{(m)}-(\log\Gamma)^{(m)}(1+v)}{m!}\right|
   &\le\frac1{m(m-1)(1-q)^m s^{m-1}}\quad(m\ge2).
 \label{eq:euler-fallback-tails}
\end{align}
To verify these inequalities subtract the locally convergent logarithmic
product. Its linear coefficient differs by
\(H_s-\log s-\gamma\in[0,1/s]\); the remaining log terms obey
\(\abs{\Log(1+v/j)-v/j}\le\abs{v/j}^2/[2(1-q)]\).
Differentiate the tail and use
\(\sum_{j>s}j^{-m}\le s^{1-m}/(m-1)\).
The finite derivative formulas are
\[
 F_s'=\log s-\sum_{j=1}^s(j+v)^{-1},\qquad
 F_s^{(m)}/m!=\frac{(-1)^m}{m}\sum_{j=1}^s(j+v)^{-m}.
\]
All denominators have real part at least \(1/4\).
A common cutoff
\(s\ge\max(2B_\Gamma,2,2(1+B_\Gamma+B_\Gamma^2)/\delta)\)
makes every truncation error at most \(\delta\).

For desired output error \(\eta\), choose \(\delta=\eta/4\).
Deliver each value logarithm to
\(\eta/[8(s+1)(B_\Gamma+1)]\), \(\log s\) to
\(\eta/[8(B_\Gamma+1)]\), each normalized reciprocal-power summand
to \(\eta/[8(s+1)]\), and the derivative's \(\log s\) to
\(\eta/8\). Exact rational summation then leaves output slack.
This takes \(O((1+B_\Gamma^2)/\delta)\) terms. Using only
output-level unit roundoff in this much longer sum need not meet the
accumulation guard. Termwise precision or exact summation is essential;
\cref{prop:gamma-provider} subsequently improves the term count.

\section{Explicit conditioning for Gauss data}
\label{app:gauss-conditioning}
For \(s\ge1\) let \(A_s=2^ss!P_s\in\mathbb Z[x]\), with leading
coefficient \(a_s=(2s)!/s!\), and let \(c_s\) be that of \(P_s\).
Define
\begin{align*}
 \sigma_s&=[a_s^{2s-1}2^{s(s-1)}]^{-1},&
 d_s&=c_s\sigma_s^{s-1},\\
 S_1&=\sum_jj\abs{[x^j]P_s},&
 S_2&=\sum_jj(j-1)\abs{[x^j]P_s},\\
 J_*&=\sum_j\abs{[x^j]J_s},&
 J_1&=\sum_jj\abs{[x^j]J_s},\\
 L_w&=2J_1/d_s+4J_*S_2/d_s^2,&
 \mu_s&=[4\max(1,2sS_1/d_s)]^{-1}.
\end{align*}
For \(s\ge2\), the nonzero integer resultant of \(A_s,A_s'\)
equals, up to sign, \(a_s^{2s-1}\) times the squared-gap product.
All gaps are at most two, so every gap is at least \(\sigma_s\).
For \(s=1\) there is no gap condition. Likewise the nonzero integers
\(A_s(\pm1)\) give endpoint distances at least
\(1/(a_s2^{s-1})\). At a root \(\abs{P_s'}\ge d_s\).
On an interval of width at most \(d_s/[2\max(1,S_2)]\), its
derivative has modulus at least \(d_s/2\); differentiation of
\(J_s/P_s'\) gives Lipschitz bound \(L_w\).

For the nodal Lagrange polynomial,
\(\abs{l_i'}\le sS_1/d_s\) and \(l_i(\xi_i)=1\).
Integrating its square toward the farther endpoint gives
\(\omega_i\ge\mu_s\). Thus it suffices to isolate to width
\begin{equation}
 \delta_*\le\min\left\{\frac18,\frac{\sigma_s}4,
 \frac{d_s}{4\max(1,S_2)},\eta_x,
 \frac{\eta_w}{2s\max(1,L_w)},
 \frac{\mu_s}{2\max(1,L_w)}\right\}.
 \label{eq:gauss-isolation-target}
\end{equation}
Take \(2^{1-h}\le\delta_*\). All conditioning lengths are
\(O(s^3\log(s+2))\), giving the depth in \cref{prop:gauss-provider}.

For weight arithmetic form
\[
 Q_s(c)=s!\int_{-1}^1\frac{A_s(t)-A_s(c)}{t-c}\dd t,
 \qquad
 [c^r]Q_s=2\sum_{\substack{j>r\\j-r\ {
m odd}}}
                   [t^j]A_s\frac{s!}{j-r}.
\]
Both integer polynomials have \(O(s\log(s+2))\)-bit coefficients.
The rational weight at a midpoint is
\(Q_s(\widehat\xi)/(s!A_s'(\widehat\xi))\).
At \(\widehat\xi=c/q\), homogeneous evaluation cancels the common
\(q^{s-1}\), giving \(O(sU_s)\)-bit numerators and denominators.
This proves the weight-size assertion used in the total cost ledger.

The global depth is sufficient, not mandatory. An isolating interval
with a certified nonzero \(P_s'\) can instead use direct interval
evaluation of \(J_s/P_s'\). If its node and true-weight enclosures
already meet the assigned radii and positivity tests, it is a valid
earlier stopping certificate for the same algorithm. This observation
does not improve the stated worst-case complexity; its bounded software
use is documented in the supporting note \cite{rec12}.
